\subsubsection{Theorem~\ref{thm:dirichlet-continuous-data-existence} 的证明}\label{subsec:perron-existence-proof}

现在已准备好证明主要结果. 设 $S_g$ 为 Dirichlet 问题的所有下解组成的集合, 并令
\begin{equation}\label{eq:perron-supremum-of-subsolutions}
u(\bm{x})=\sup_{v\in S_g}v(\bm{x}).
\end{equation}
由于每个下解都小于或等于每个上解, 且下解和上解都存在, \eqref{eq:perron-supremum-of-subsolutions} 中的上确界有限, 因而 $u$ 良定义.

\begin{Lemma}
\label{lem:perron-envelope-harmonic}
函数 $u$ 在 $\Omega$ 内调和.
\end{Lemma}

\begin{Proof}
取 $\bm{x}$ 为 $\Omega$ 内任意一点, 并取 $S_g$ 中的序列 $v_m$, 使 $v_m(\bm{x})\to u(\bm{x})$. 显然 $v_m$ 有上界 (任意上解都是上界), 也可假设它有下界; 必要时用 $\max(v_m,v_0)$ 替换 $v_m$, 其中 $v_0$ 是任意下解. 选取 $R$, 使 $B=B_R(\bm{x})$ 的闭包包含于 $\Omega$ 内, 设 $V_m$ 为 $v_m$ 关于 $B$ 的调和提升. 那么 $V_m(\bm{x})\to u(\bm{x})$, 且由 Theorem~\ref{thm:harmonic-sequence-compactness}, $V_m$ 有子列 $V_{m_k}$, 在 $B$ 上收敛到调和函数 $v$; 收敛在 $B$ 的每个紧子集上都是一致的. 显然, $B$ 内有 $v\leq u$. 现在证明事实上 $B$ 内 $v=u$. 这显然蕴含引理.

反设对某个 $\bm{y}\in B$ 有 $v(\bm{y})<u(\bm{y})$. 那么存在函数 $W\in S_g$, 使 $v(\bm{y})<W(\bm{y})$. 令 $w_k=\max(W,V_{m_k})$, 并设 $W_k$ 为 $w_k$ 关于 $B$ 的调和提升. 与前面一样, $W_k$ 的一个子列收敛到 $B$ 内的调和函数 $w$. 显然在 $B$ 上 $v\leq w$, 且 $v(\bm{x})=w(\bm{x})$. 由最大值原理, 这蕴含在 $B$ 上 $v=w$, 与 $W$ 的选取矛盾.
\end{Proof}

剩下需要证明 $u$ 实际上直到 $\partial\Omega$ 都连续, 并取到给定的边界数据. 设 $\bm{\xi}$ 为 $\partial\Omega$ 上一点. 由于假设了 $\partial\Omega$ 属于 $C^2$ 类, 存在位于 $\Omega$ 外部的球 $B=B_R(\bm{y})$, 使 $\overline B\cap\overline\Omega=\{\bm{\xi}\}$. 现在定义
\begin{equation}\label{eq:perron-exterior-ball-barrier}
w(\bm{x})=\begin{cases}R^{2-n}-|\bm{x}-\bm{y}|^{2-n},&n\geq3\\\log\dfrac{|\bm{x}-\bm{y}|}{R},&n=2.\end{cases}
\end{equation}
那么
\begin{enumerate}
\item $w$ 在 $\Omega$ 内超调和 (实际上调和).
\item 在 $\overline\Omega\setminus\{\bm{\xi}\}$ 上 $w>0$, 且 $w(\bm{\xi})=0$.
\end{enumerate}

具有性质 (1) 和 (2) 的连续函数称为 $\bm{\xi}$ 处关于 $\Omega$ 的\emph{屏障}. 有以下结果.

\begin{Lemma}
\label{lem:perron-boundary-continuity}
设 $u$ 为上述构造的 $\Omega$ 内的调和函数, 并设 $\bm{\xi}\in\partial\Omega$. 那么当 $\bm{x}\to\bm{\xi}$ 时, $u(\bm{x})\to g(\bm{\xi})$.
\end{Lemma}

\begin{Proof}
\PDEIndexMark{idx:perron-boundary-continuity-proof}选取 $\epsilon>0$, 令 $M=\max_{\partial\Omega}|g|$. 设 $w$ 为 $\bm{\xi}$ 处的屏障, 取 $\delta$ 和 $k$, 使对 $\bm{x}\in\partial\Omega$, $|\bm{x}-\bm{\xi}|<\delta$, 有 $|g(\bm{x})-g(\bm{\xi})|<\epsilon$, 并对 $\bm{x}\in\Omega$, $|\bm{x}-\bm{\xi}|\geq\delta$, 有 $kw(\bm{x})\geq2M$. 函数 $g(\bm{\xi})+\epsilon+kw(\bm{x})$ 与 $g(\bm{\xi})-\epsilon-kw(\bm{x})$ 分别为上解和下解, 因而
\begin{equation}\label{eq:perron-barrier-boundary-comparison}
g(\bm{\xi})-\epsilon-kw(\bm{x})\leq u(\bm{x})\leq g(\bm{\xi})+\epsilon+kw(\bm{x}).
\end{equation}
由于当 $\bm{x}\to\bm{\xi}$ 时 $w(\bm{x})\to0$, 引理立即成立.
\end{Proof}

由证明可见, 只要能够构造屏障, 就可以放宽 $\partial\Omega$ 的光滑性假设. 上述构造对所有凸区域都适用. 在其他情形下, 也常可找到其他屏障, 例如见 Problem~\ref{ex:planar-local-logarithmic-barrier}.

\paragraph*{Problems}

\begin{enumerate}[label=\thesection.\arabic*.,ref=\thesection.\arabic*,start=8,series=book-problems]

\item\label{ex:poisson-formula-harmonicity} 验证 \eqref{eq:poisson-ball-solution-formula} 给出的函数是调和的.

\item\label{ex:subharmonic-harmonic-lifting} 证明次调和函数小节中的第二条结论.

\item\label{ex:smooth-subharmonic-laplacian} 证明, 若 $u$ 属于 $C^2$ 类, 且在本节所定义的意义下次调和, 则 $\Delta u\geq0$.

\item\label{ex:compact-equicontinuity-uniform} 设序列 $f_m$ 在紧集 $S$ 的每一点处等度连续. 证明它在 $S$ 上一致等度连续.

\item\label{ex:punctured-disk-dirichlet-nonexistence} 令 $\Omega=\{(x,y)\in\mathbb{R}^2\mid0<x^2+y^2<1\}$. 证明 Dirichlet 问题在 $\Omega$ 内 $\Delta u=0$, 当 $x^2+y^2=1$ 时 $u(\bm{x})=1$, 且 $u(0)=0$, 无解. Hint: 先证明若存在解, 则也存在径向对称解.

\item\label{ex:local-barrier-to-global-barrier} 若定义中的条件 (i) 和 (ii) 只在 $\bm{\xi}$ 的某个邻域 $N$ 内成立, 则称函数 $w$ 为 $\bm{\xi}$ 处关于 $\Omega$ 的局部屏障. 证明若存在局部屏障, 则存在屏障. Hint: 取以 $\bm{\xi}$ 为球心的小球 $B$, 并令 $m=\inf_{N\setminus B}w$. 在 $B$ 上定义 $\bar w=\min(m,w)$, 在其他地方定义 $\bar w=m$.

\item\label{ex:planar-local-logarithmic-barrier} 设 $\Omega$ 为 $\mathbb{R}^2$ 中的区域, 原点在其边界上. 设 $r$ 和 $\theta$ 表示极坐标. 假设在 $\Omega\cap N$ 内存在 $\theta$ 的单值分支, 其中 $N$ 为原点的一个邻域. 证明原点处存在局部屏障. Hint: 考虑解析函数 $1/\log z$.

\item\label{ex:perron-proof-unbounded-domain} 为什么 Theorem~\ref{thm:dirichlet-continuous-data-existence} 的证明不适用于无界区域?

\item\label{ex:unit-ball-volume-gamma-function} 证明 $\omega_n=\pi^{n/2}/\Gamma(\frac n2+1)$. Hint: 分别在 Cartesian 坐标和极坐标中计算 \allowbreak $\int_{\mathbb{R}^n}\exp(-|\bm{x}|^2)\,\mathrm{d}\bm{x}$. 比较这两个表达式, 得到单位球面的面积 $\Omega_n=n\omega_n$ 的一个公式.

\end{enumerate}
